%! Introduction to Eigenvalues — Unit 6, Lesson 6.1 — Warm-Up (Answer Key)
\documentclass[10pt]{article}
\usepackage{linalg-article}
\usepackage{linalg-key}   % pulls in -boxes; do NOT also load -boxes
\newcommand{\TallMath}[1]{$\displaystyle #1\rule[-1.4em]{0pt}{3.2em}$}

\begin{document}

\pageheader{Unit 6, Lesson 6.1}{Warm-Up --- Answer Key}
\namedateperiod

\vspace{0.12in}

\begin{notesbox}{Getting Ready: The Three Moves We'll Use Today}
\small Yesterday you \emph{tested} vectors you were handed. Today we \emph{find} the special
directions ourselves --- and these three moves are the whole method. Answer quickly.

\vspace{4pt}
\textbf{1. Subtract $\lambda$ down the diagonal.} For $A=\begin{bmatrix} 4 & 1 \\ 2 & 3 \end{bmatrix}$,
write the matrix $A-\lambda I$ (subtract the unknown $\lambda$ from each diagonal entry):
\[
  A-\lambda I=\begin{bmatrix}{\color{keyred}\mathbf{4-\lambda}} & {\color{keyred}\mathbf{1}}\\[3pt]{\color{keyred}\mathbf{2}} & {\color{keyred}\mathbf{3-\lambda}}\end{bmatrix}.
\]

\vspace{4pt}
\textbf{2. The $2\times2$ determinant (Unit 5).} Recall $\det\begin{bsmallmatrix} a & b \\ c & d
\end{bsmallmatrix}=ad-bc$. Compute
\[
  \det\begin{bmatrix} 4 & 1 \\ 2 & 3 \end{bmatrix} = {\color{keyred}\mathbf{(4)(3)-(1)(2)=10}}.
\]

\vspace{4pt}
\textbf{3. Solve a quadratic (factor).} Find both values of $\lambda$ that make this true:
\[
  \lambda^2-5\lambda+6=0 \qquad\Longrightarrow\qquad \lambda={\color{keyred}\mathbf{2}}\ \text{ or }\ \lambda={\color{keyred}\mathbf{3}}.
\]
\end{notesbox}

\begin{teachernote}
These three moves \emph{are} today's method for finding eigenvalues, in order:
\textbf{(1)} form $A-\lambda I$ (subtract $\lambda$ from the diagonal), \textbf{(2)} take its
$2\times2$ determinant, and \textbf{(3)} set that determinant to $0$ and solve the quadratic ---
its roots are the eigenvalues. Item~1 uses the same $A=\begin{bsmallmatrix}4&1\\2&3\end{bsmallmatrix}$
students will meet again in the homework; item~2 is that determinant with $\lambda=0$; item~3 rehearses
factoring $\lambda^2-5\lambda+6=(\lambda-2)(\lambda-3)$. Debrief all three quickly, then name the plan:
``today we do exactly these three moves \emph{with a $\lambda$ inside.}''
\end{teachernote}

\end{document}
